\documentclass[11pt,a4paper]{article}
\usepackage[T1]{fontenc}
\usepackage{lmodern,amsmath,amssymb,amsthm,microtype}
\usepackage[margin=27mm]{geometry}
\usepackage[colorlinks=true,linkcolor=blue,urlcolor=blue]{hyperref}
\newtheorem{proposition}{Proposition}
\newcommand{\Q}{\mathbb Q}\newcommand{\Z}{\mathbb Z}
\setlength{\emergencystretch}{2em}
\title{A two-label monic Lucas clause}
\author{}\date{17 November 2026}
\begin{document}\maketitle

\begin{proposition}
Let $s,b\in\Z$, with $|s|\ge3$ and $b\ne0$. Put
$a=s^2$ and $d=a-2$. The following are equivalent:
\begin{enumerate}
\item There is a rational number $z$ satisfying the monic equation
\begin{equation}\label{eq:monic}
 z^2=(a^2-4)d^2b^2+4.
\end{equation}
\item There is an integer $t$ such that
\begin{equation}\label{eq:lucas}
 t^2-a(a-2)tb+(a-2)^2b^2=1.
\end{equation}
\end{enumerate}
Every rational root $z$ in the first assertion is integral, and
\[
 t=\frac{z+adb}{2}
\]
is an integer satisfying the second. Every such integer satisfies
\[
 \log|t|>s^2-3.
\]
\end{proposition}
\begin{proof}
The radicand in \eqref{eq:monic} is an integer, so a rational square
root is integral. Furthermore
\[
 z^2-(adb)^2=4(1-d^2b^2).
\]
Squares modulo four show that $z$ and $adb$ have the same parity;
therefore the displayed $t$ is integral. Substitution gives
\eqref{eq:lucas}. Conversely $z=2t-adb$ gives
\eqref{eq:monic}. These assertions do not require $s$ or $b$ to be
odd. The growth estimate is the elementary Lucas lemma of
\cite{Lucas}, applied to $a=s^2>7$.
\end{proof}

The radicand is positive because $a\ge9$ and $b\ne0$.
The two roots are distinct and nonzero. Thus, when the two integer
labels $s,b$ are available as exact rational elements of a source
algebra, the growth clause needs only one additional finite quadratic
root and no additional free rational coordinate. In the existing
elliptic source the functions $L_S,L_Z$ only approximate the integer
indices at selected primes; they are not those indices. The proposition
does not justify substituting these functions for $s,b$, and it does
not give a smaller defining-variable bound.

For completeness at the numerical-label level, put
$U_0(a)=0$, $U_1(a)=1$ and
$U_{j+1}(a)=aU_j(a)-U_{j-1}(a)$. With
\[
 k=2(a-2),\qquad b=U_k(a)/(a-2),\qquad t=U_{k+1}(a),
\]
the congruence $U_k(a)\equiv k\pmod{a-2}$ makes $b$ integral,
and the Lucas norm identity gives \eqref{eq:lucas}.
The preceding proposition then supplies $z=2t-adb$.

\begin{thebibliography}{9}\raggedright
\bibitem{Lucas}\emph{An explicit polynomial in 29 variables over the
rationals}, repository manuscript, elementary Lucas growth lemma,
\nolinkurl{papers/few_variables/few_variables_20261108_v8.tex},
lines 349--397.
\end{thebibliography}
\end{document}
